\MathBookSourcePage{229}
\subsection{化为加权估计}
\label{sec:ch08-reduction-to-weighted-estimates}
\setcounter{thmcount}{0}
\setcounter{equation}{0}

对 $z\in\overline K^z$, 令 $\delta^z\in C_0^\infty(K^z)$ 满足
\MBSourceItem{1}
\begin{equation}
\MathBookFormulaID{formula-ch08-smoothed-delta-properties}
\label{eq:ch08-smoothed-delta-properties}
\begin{gathered}
\int_\Omega\delta^z(x)P(x)\,dx=P(z)
\qquad\forall P\in\mathcal P_{k_{\max}},\\
\|\nabla_k\delta^z\|_{L^\infty(\Omega)}\leq C_kh^{-d-k},
\qquad k=0,1,\ldots,
\end{gathered}
\end{equation}
其中常数 $C_k$ 只依赖于剖分, $\nabla_k$ 表示所有 $k$ 阶导数组成的向量(参见习题~\ref{exr:ch08-07}). 根据构造,
\MBSourceItem{2}
\begin{equation}
\MathBookFormulaID{formula-ch08-smoothed-delta-gradient-reproduction}
\label{eq:ch08-smoothed-delta-gradient-reproduction}
\nu\cdot\nabla v(z)=(\nu\cdot\nabla v,\delta^z)
\qquad\forall v\in V_h
\end{equation}
对任意方向向量 $\nu$ 成立.

定义 $g^z\in V$ 为伴随变分问题
\[
\MathBookFormulaID{formula-ch08-weighted-adjoint-green-problem}
a(v,g^z)=(-\nu\cdot\nabla\delta^z,v)
\qquad\forall v\in V
\]
的解. 注意 $\delta^z$ 和 $g^z$ 都依赖于 $h$, 但为简化记号省略这一点. 令 $g_h^z$ 表示 $g^z$ 的有限元逼近:
\[
\MathBookFormulaID{formula-ch08-weighted-adjoint-green-galerkin}
a(v,g_h^z)=(-\nu\cdot\nabla\delta^z,v)
\qquad\forall v\in V_h.
\]
由 $u,u_h,g^z,g_h^z$ 的定义和分部积分,
\MBSourceItem{3}
\begin{align}
\MathBookFormulaID{formula-ch08-localization-identity}
\nu\cdot\nabla u_h(z)
&=(\nu\cdot\nabla u_h,\delta^z)=(u_h,-\nu\cdot\nabla\delta^z)
=a(u_h,g_h^z)=a(u_h,g_h^z-g^z)+a(u_h,g^z)\notag\\
&=a(u,g^z)-a(u,g^z-g_h^z)
=(-\nu\cdot\nabla\delta^z,u)-a(u,g^z-g_h^z)\notag\\
&=(\delta^z,\nu\cdot\nabla u)-a(u,g^z-g_h^z).
\label{eq:ch08-localization-identity}
\end{align}
由 Hölder 不等式,
\[
\MathBookFormulaID{formula-ch08-adjoint-error-duality-bound}
|a(g^z-g_h^z,u)|\leq C\|g^z-g_h^z\|_{W_1^1(\Omega)}\|u\|_{W_\infty^1(\Omega)}
\]
\[
\MathBookFormulaID{formula-ch08-adjoint-error-weighted-bound}
\leq C\left(\int_\Omega\sigma_z^{d+\lambda}|\nabla(g^z-g_h^z)|^2\,dx\right)^{1/2}
\left(\int_\Omega\sigma_z^{-d-\lambda}\,dx\right)^{1/2}\|u\|_{W_\infty^1(\Omega)}.
\]
令
\MBSourceItem{4}
\begin{equation}
\MathBookFormulaID{formula-ch08-weighted-green-error-supremum}
\label{eq:ch08-weighted-green-error-supremum}
M_h:=\sup_{z\in\overline\Omega}
\left(\int_\Omega\sigma_z^{d+\lambda}
\bigl(|\nabla(g^z-g_h^z)|^2+(g^z-g_h^z)^2\bigr)\,dx\right)^{1/2},
\end{equation}
则
\MBSourceItem{5}
\begin{equation}
\MathBookFormulaID{formula-ch08-green-error-action-bound}
\label{eq:ch08-green-error-action-bound}
|a(g^z-g_h^z,u)|\leq CM_h
\left(\int_\Omega\sigma_z^{-d-\lambda}\,dx\right)^{1/2}\|u\|_{W_\infty^1(\Omega)}
\leq CM_h\lambda^{-1/2}(\kappa h)^{-\lambda/2}\|u\|_{W_\infty^1(\Omega)}.
\end{equation}
因此, 只需证明下述结论.

\MBSourceItem{6}
\begin{Lemma}
\label{lem:ch08-weighted-green-error}
适当选取 $\lambda>0$, 并令 $\kappa$ 足够大, 则存在 $h_0>0$, 使得对所有 $0<h<h_0$,
\[
\MathBookFormulaID{formula-ch08-weighted-green-error-lemma}
\sup_{z\in\overline\Omega}\int_\Omega\sigma_z^{d+\lambda}
\bigl(|\nabla(g^z-g_h^z)|^2+(g^z-g_h^z)^2\bigr)\,dx\leq Ch^\lambda.
\]
\end{Lemma}

下一节将通过一系列命题给出证明. 作为推论, 立即得到下述结果.

\MBSourceItem{7}
\begin{Corollary}
\label{cor:ch08-weighted-green-stability}
存在 $C<\infty$ 和 $h_0>0$, 使得对所有 $0<h<h_0$,
\[
\MathBookFormulaID{formula-ch08-weighted-green-stability}
\|g^z-g_h^z\|_{W_1^1(\Omega)}\leq C.
\]
\end{Corollary}

由式~\eqref{eq:ch08-localization-identity}, 还有
\[
\MathBookFormulaID{formula-ch08-localization-functional}
(\delta^z,\nu\cdot\nabla u)-\nu\cdot\nabla u_h(z)=a(u-v,g^z-g_h^z)
\qquad\forall v\in V_h.
\]
应用引理~\ref{lem:ch08-weighted-green-error}, 可证明如下局部化结果.

\MBSourceItem{8}
\begin{Corollary}
\label{cor:ch08-localized-gradient-average}
对引理~\ref{lem:ch08-weighted-green-error} 中的 $\lambda>0$, 存在 $C<\infty$ 和 $h_0>0$, 使得对所有 $0<h<h_0$ 以及所有 $v\in V_h$,
\[
\MathBookFormulaID{formula-ch08-localized-gradient-average}
|(\delta^z,\nu\cdot\nabla u)-\nu\cdot\nabla u_h(z)|
\leq C\inf_{v\in V_h}
\left(\int_\Omega h^\lambda\sigma_z^{-d-\lambda}|\nabla(u-v)|^2\,dx\right)^{1/2}.
\]
\end{Corollary}

这个结果表明, $\nu\cdot\nabla u_h(z)$ 非常接近 $z$ 附近 $\nu\cdot\nabla u$ 的某种平均值, 而其差由逼近误差(即最佳逼近误差)控制. 这里的误差是局部化的, 因为权函数 $h^\lambda\sigma_z^{-d-\lambda}(x)$ 是一种光滑化的 Dirac $\delta$ 函数, 当 $x$ 远离 $z$ 时衰减到零, 至少与 $|x-z|^{-d}$ 一样快.
